% ch01-clean-architecture — المعماريّة النظيفة (Clean Architecture)
\section{المعماريّة النظيفة (Clean Architecture)}

في الفصل السابق اتّفقنا على المصطلحات الأساسية: الخلفية \en{(backend)}، والواجهة
\en{(frontend)}، وقاعدة البيانات، والـ \en{API}، وفكرة «الطبقة» بشكل عام. في هذا
الفصل نفتح صندوق الخلفية ونشوف كيف اتبنى من الداخل. الخلفية في \en{ScholarPath}
مكتوبة بـ \en{.NET 10}، ومنظَّمة بأسلوب اسمه \en{Clean Architecture} — أو
«المعماريّة النظيفة». وده أهم قرار هندسيّ في المشروع كلّه، فخلّينا نفهمه واحدة واحدة.

الفكرة ببساطة: بدل ما نكتب الكود كلّه مكوَّمًا في مكان واحد، نقسّمه إلى \textbf{أربع
طبقات متداخلة} زيّ حلقات دائرة. كل طبقة لها مسؤولية واحدة واضحة، والأهمّ: قاعدة صارمة
واحدة تحكم الاعتماد بينهم.

\begin{idea}[الفكرة: الاعتماد يروح للداخل فقط]
تخيّل أربع دوائر متداخلة. في القلب توجد \textbf{قواعد العمل} (منطق المِنَح
والحجوزات والدفع)، وفي الخارج توجد \textbf{التفاصيل التقنية} (قاعدة البيانات،
\en{Azure}، \en{Stripe}). القاعدة الذهبية: \textbf{الاعتماد يتّجه للداخل دائمًا}.
الطبقة الداخلية لا تعرف شيئًا عن الطبقة الخارجية. القلب (منطق العمل) لا يعرف أنّ
قاعدة البيانات اسمها \en{SQL Server}، ولا أنّ مزوّد الذكاء الاصطناعي اسمه
\en{Azure OpenAI}. هو يعرف \emph{ماذا} يحتاج، لا \emph{كيف} يُنفَّذ.
\end{idea}

\subsection{الطبقات الأربع ومسؤولية كل واحدة}

النظام مبنيّ على أربعة مشاريع \en{(.csproj)} منفصلة، كل واحد منها طبقة. ترتيبها من
الداخل للخارج هكذا: \en{Domain} ثم \en{Application} ثم \en{Infrastructure} ثم
\en{API}. اتجاه الاعتماد يمشي في العكس — من الخارج للداخل:

\begin{center}\begin{LTR}
\en{API} $\longrightarrow$ \en{Infrastructure} $\longrightarrow$ \en{Application} $\longrightarrow$ \en{Domain}
\end{LTR}\end{center}

يعني \en{API} يعتمد على مَن تحته، و\en{Domain} في القلب لا يعتمد على أحد إطلاقًا.
ده مضمون فعليًّا في الكود عن طريق مراجع المشاريع \en{(project references)}: مشروع
\code{ScholarPath.Domain} لا يحمل أي مرجع لأي مشروع آخر، ومشروع
\code{ScholarPath.Application} يشير إلى \code{ScholarPath.Domain} فقط. المُصرِّف
\en{(compiler)} نفسه يرفض أي محاولة لكسر هذا الترتيب.

\begin{center}
\begin{tabular}{L{2.3cm} L{3.7cm} L{8.8cm}}
\toprule
\textbf{\color{spPrimary}الطبقة} & \textbf{\color{spPrimary}المشروع \en{(csproj)}} & \textbf{\color{spPrimary}المحتوى الأساسي} \\
\midrule
\en{Domain} & \code{ScholarPath.Domain} & الكيانات \en{(entities)} $\approx$55، والـ \en{enums}، وأحداث الدومين، والمنفذان \code{ICurrentUserService} و\code{IDateTimeService}. بلا أي اعتماد على \en{framework}. \\
\en{Application} & \code{ScholarPath.Application} & حالات الاستخدام كأوامر واستعلامات \en{CQRS} عبر \en{MediatR}، والـ \en{validators}، وقوالب \en{AutoMapper}، و$\approx$33 \en{port} (واجهة) يجب على العالم الخارجي تنفيذها. \\
\en{Infrastructure} & \code{ScholarPath.Infrastructure} & $\approx$46 \en{adapter} تنفّذ الـ \en{ports}: \en{EF Core}، \en{Stripe}، \en{Azure OpenAI}، \en{ACS}، \en{Blob}، \en{Key Vault}، \en{SignalR}، \en{Hangfire}. \\
\en{API} & \code{ScholarPath.API} & الـ \en{controllers} الرفيعة، ومستقبِل \en{Stripe webhook}، والهُويّة \en{(ASP.NET Identity + JWT)}. \\
\bottomrule
\end{tabular}
\end{center}

خلّينا ندخل على كل طبقة بالتفصيل، من القلب للخارج.

\subsection{طبقة \en{Domain} — القلب النقيّ}

هذه أعمق طبقة، وأهم قاعدة فيها: \textbf{صفر اعتماد على أي \en{framework}}. لا تعرف
\en{Entity Framework}، ولا \en{ASP.NET}، ولا \en{Azure}. هي مجرّد \en{C\#} خالص
يصف عالم المشكلة: كيان الطالب، كيان المِنحة، كيان الحجز، حالات الدفع كـ \en{enum}،
وأحداث الدومين \en{(domain events)} التي تُطلَق عند حدوث شيء مهمّ.

الاستثناء الوحيد أنّ الدومين يُعرِّف \textbf{منفذين} \en{(ports)} فقط، لأنّهما مفهومان
ينتميان لقلب العمل نفسه: مَن هو المستخدم الحاليّ؟ وكم الساعة الآن؟ خلّينا نشوفهما.

هنشوف إيه: الملف الحقيقيّ الذي يُعرِّف المنفذين داخل الدومين.

\begin{english}
\begin{lstlisting}[language={[Sharp]C}]
namespace ScholarPath.Domain.Interfaces;

public interface ICurrentUserService
{
    Guid? UserId { get; }
    string? Email { get; }
    string? ActiveRole { get; }
    IReadOnlyCollection<string> Roles { get; }
    bool IsAuthenticated { get; }
    string? IpAddress { get; }
    // ...
}

public interface IDateTimeService
{
    DateTimeOffset UtcNow { get; }
    DateOnly Today { get; }
}
\end{lstlisting}
\end{english}

الشرح سطرًا بسطر لأهم الأسطر:
\begin{itemize}
\item \code{ICurrentUserService}: منفذ يجيب على سؤال «مَن يُنفّذ هذا الطلب الآن؟» —
  مُعرِّفه \code{UserId}، بريده، ودوره النشِط \code{ActiveRole}. الدومين يحتاج هذه
  المعلومة لتطبيق قواعد الملكية (مثلًا: هل هذا الطالب يملك هذا الطلب؟)، لكنّه لا
  يعرف \emph{من أين} تأتي — من \en{JWT}؟ من \en{cookie}؟ لا يهمّه.
\item \code{IDateTimeService}: منفذ بسيط يعطي \code{UtcNow} و\code{Today}. لماذا لا
  نستدعي \code{DateTime.Now} مباشرة؟ لأنّ تجميد الوقت خلف منفذ يجعل الكود
  \textbf{قابلًا للاختبار}: في الاختبار نُثبّت «الآن» على لحظة معيّنة ونتأكّد أنّ
  منطق المواعيد النهائية يتصرّف صحّ.
\end{itemize}

\begin{why}
لماذا نُبقي الدومين نقيًّا بلا \en{framework}؟ لأنّ قواعد العمل هي أثمن ما نملك،
وأطول ما يعيش. الأُطُر التقنية \en{(frameworks)} تتغيّر: قد نستبدل \en{SQL Server}
بغيره، أو نُرقّي \en{.NET}. لو كانت قواعد العمل مشتبكة مع التفاصيل التقنية، فأي
تغيير تقنيّ يهدّد المنطق. عزلُها يعني أنّ منطق المِنَح يبقى ثابتًا مهما تبدّل ما
حوله — وده يجعله قابلًا للاختبار بلا قاعدة بيانات ولا شبكة.
\end{why}

\subsection{طبقة \en{Application} — حالات الاستخدام والمنافذ}

فوق الدومين تأتي طبقة \en{Application}. مسؤوليتها: تنسيق \textbf{حالات الاستخدام}
\en{(use cases)} — «الطالب يحجز جلسة»، «المستشار يقبل الحجز»، «أرشد لي منَحًا». كل
حالة استخدام مكتوبة على نمط \en{CQRS}.

\begin{keyp}[نقطة مهمة: \en{CQRS}]
\en{CQRS} اختصار لـ \en{Command Query Responsibility Segregation}. الفكرة: نفصل كل
عملية إلى نوعين. \textbf{أمر} \en{(Command)} يُغيّر الحالة (كتابة) مثل
\code{AcceptBookingCommand}. و\textbf{استعلام} \en{(Query)} يقرأ فقط بلا أي أثر
جانبيّ مثل \code{GetMyRecommendationsQuery}. كل أمر أو استعلام له مُعالِج
\en{(handler)} واحد مسؤول عنه، فيصبح كل مِلفّ صغيرًا ومملوكًا لمطوّر واحد.
\end{keyp}

هذه الطبقة تحتوي أيضًا على الـ \en{validators} (تتحقّق من صحّة المُدخَلات قبل أي
عمل) عبر \en{FluentValidation}، وقوالب \en{AutoMapper} (تحوّل الكيانات إلى
\en{DTOs} للعرض). لكن الجزء الأهمّ معماريًّا: هنا تُعرَّف الـ \textbf{$\approx$33
منفذًا} \en{(ports)}. المنفذ هو مجرّد \en{interface} — وعدٌ بخدمة يحتاجها المنطق دون
أن يعرف تنفيذها. مثال حيّ: منفذ الذكاء الاصطناعي.

هنشوف إيه: منفذ \code{IAiService} الحقيقيّ — عقد الذكاء الاصطناعي كما تراه طبقة العمل.

\begin{english}
\begin{lstlisting}[language={[Sharp]C}]
namespace ScholarPath.Application.Common.Interfaces;

public interface IAiService
{
    Task<AiRecommendationResult> GenerateRecommendationsAsync(
        Guid userId, int topN, CancellationToken ct);

    Task<AiEligibilityResult> CheckEligibilityAsync(
        Guid userId, Guid scholarshipId, CancellationToken ct);

    Task<AiChatResponse> AskAsync(
        Guid userId, string sessionId, string message,
        IReadOnlyList<AiChatHistoryTurn> history, CancellationToken ct);
}
\end{lstlisting}
\end{english}

الشرح: هذا الـ \en{interface} يصف \emph{ماذا} يقدّم الذكاء الاصطناعي في ثلاث
عمليات فقط: \code{GenerateRecommendationsAsync} للترشيحات،
و\code{CheckEligibilityAsync} لفحص الأهليّة، و\code{AskAsync} للمحادثة. لاحِظ أنّه
\textbf{لا يوجد ذِكر لـ \en{Azure} ولا \en{OpenAI} ولا أي مِفتاح}. طبقة العمل تعرف
أنّها تحتاج «خدمة ذكاء اصطناعي» بهذا الشكل، ولا تهتمّ من يقف خلفها. هذا هو معنى
«المنفذ»: عقدٌ نظيف يعيش في الداخل.

\subsection{طبقة \en{Infrastructure} — المحوّلات التي تُنفّذ الوعود}

كل منفذ في طبقة \en{Application} لا بدّ أن يجد مَن يُنفّذه في الواقع. هنا يأتي دور
طبقة \en{Infrastructure}: هي بيت الـ \textbf{$\approx$46 محوّلًا} \en{(adapters)}.
المحوّل هو الصف \en{(class)} الذي يقول «أنا أُنفّذ هذا المنفذ باستخدام تقنية
حقيقية». هنا يعيش \en{EF Core} (قاعدة البيانات)، و\en{Stripe} (الدفع)، و\en{Azure
OpenAI} (الذكاء الاصطناعي)، و\en{ACS} (الفيديو)، و\en{Blob} (تخزين الملفات)،
و\en{Key Vault} (المفاتيح السرّية)، و\en{SignalR} (الوقت الحقيقيّ)، و\en{Hangfire}
(المهام الخلفية).

هنشوف إيه: محوّل \code{AzureOpenAiService} الحقيقيّ الذي يُنفّذ منفذ \code{IAiService}
الذي رأيناه للتوّ.

\begin{english}
\begin{lstlisting}[language={[Sharp]C}]
public sealed class AzureOpenAiService(
    IHttpClientFactory httpFactory,
    IOptions<AiOptions> opts,
    IKnowledgeRetriever retriever,
    LocalAiService local,
    ILogger<AzureOpenAiService> logger) : IAiService
{
    // Recommendations & eligibility are deterministic -> delegate to the local engine
    public Task<AiRecommendationResult> GenerateRecommendationsAsync(
        Guid userId, int topN, CancellationToken ct)
        => local.GenerateRecommendationsAsync(userId, topN, ct);

    public async Task<AiChatResponse> AskAsync(...)
    {
        // If Azure endpoint/key are not set, answer with the local RAG router
        if (!IsAzureChatConfigured(ai))
            return await local.AskAsync(userId, sessionId, msg, history, ct);
        // ... retrieve (RAG) -> augment -> call Azure OpenAI -> on failure fall back to local
    }
}
\end{lstlisting}
\end{english}

الشرح لأهمّ الأسطر:
\begin{itemize}
\item السطر \code{... : IAiService}: إعلان صريح بأنّ هذا الصفّ \textbf{يُنفّذ}
  المنفذ. المُصرِّف يُلزمه بتقديم العمليات الثلاث كلّها.
\item السطر \code{=> local.GenerateRecommendationsAsync(...)}: الترشيحات ليست نداءً
  لنموذج لغويّ، بل حساب حتميّ. لذلك يُفوّض المحوّل السحابيّ هذه العملية لمحوّل محلّي
  \code{LocalAiService} — نفس المنفذ، تنفيذ مختلف حسب الحاجة.
\item كتلة الـ \code{if} والتعليقات: لو لم يُضبَط مِفتاح \en{Azure}، أو فشل النداء
  الشبكيّ، يتراجع المحوّل بأمان إلى المسار المحلّي بدل أن يُفشل الطلب بخطأ 500.
  المستخدم يحصل على إجابة دائمًا.
\end{itemize}

\begin{note}
لاحِظ اتجاه الاعتماد: محوّل \code{AzureOpenAiService} في الطبقة الخارجية يعرف
المنفذ \code{IAiService} الموجود في الداخل ويُنفّذه. العكس مستحيل: طبقة
\en{Application} لا تعرف اسم \code{AzureOpenAiService} أبدًا. هذا هو «قلب الاعتماد»
\en{(Dependency Inversion)} في الممارسة.
\end{note}

\subsection{طبقة \en{API} — الواجهة الرفيعة}

أعلى طبقة هي \en{API}. مهمّتها الوحيدة: استقبال طلب \en{HTTP}، ترجمته إلى أمر أو
استعلام، وتمريره إلى \en{MediatR}، ثم إعادة النتيجة. الـ \en{controllers} هنا
\textbf{رفيعة جدًّا} (سطر أو ثلاثة) — لا منطق عمل فيها إطلاقًا. تضمّ هذه الطبقة
أيضًا مستقبِل \en{Stripe webhook} (الذي يتحقّق من توقيع \en{Stripe} ويطابق
المدفوعات)، وإعداد الهُويّة \en{(ASP.NET Identity + JWT)}، وربط كل شيء معًا في
\code{Program.cs}.

\subsection{\en{Ports \& Adapters} — أهمّ قرار في النظام}

وصلنا لجوهر الفصل. الفكرة التي تربط كل ما سبق اسمها \en{Ports \& Adapters}
(المنافذ والمحوّلات): طبقة العمل تُعلن \textbf{منافذ} (واجهات نظيفة)، وطبقة
\en{Infrastructure} تقدّم \textbf{محوّلات} تُنفّذها. اختيار أيّ محوّل يُستخدَم يحدث
عند تشغيل التطبيق في مكان واحد اسمه \code{DependencyInjection.cs}، وبناءً على
\textbf{الإعداد فقط}، بلا تغيير أي سطر كود.

هنشوف إيه: مقتطف حقيقيّ من \code{DependencyInjection.cs} يختار المحوّل حسب الإعداد.

\begin{english}
\begin{lstlisting}[language={[Sharp]C}]
// Stripe: real client only with a genuine secret key; else the dev stub.
var stripeSecretKey = config.GetValue<string>($"{StripeOptions.SectionName}:SecretKey");
if (stripeSecretKey is not null && stripeSecretKey.StartsWith("sk_"))
    services.AddSingleton<IStripeService, StripeService>();
else
    services.AddSingleton<IStripeService, StubStripeService>();

// AI provider is config-driven (Ai:Provider): AzureOpenAi | OpenAi | Local
var aiProvider = config.GetValue<string>($"{AiOptions.SectionName}:Provider") ?? "Stub";
if (string.Equals(aiProvider, "AzureOpenAi", StringComparison.OrdinalIgnoreCase))
    services.AddScoped<IAiService, AzureOpenAiService>();
else
    services.AddScoped<IAiService>(sp => sp.GetRequiredService<LocalAiService>());

// Antivirus: real ClamAV only when FileScanning:Enabled = true.
if (config.GetValue<bool>($"{FileScanningOptions.SectionName}:Enabled"))
    services.AddSingleton<IFileScanService, ClamAvFileScanService>();
else
    services.AddSingleton<IFileScanService, NoOpFileScanService>();
\end{lstlisting}
\end{english}

الشرح: كل كتلة \code{if/else} تختار محوّلًا حقيقيًّا أو بديلًا للتطوير
\en{(stub)} بناءً على قيمة في الإعداد:
\begin{itemize}
\item مِفتاح \code{Stripe:SecretKey} يبدأ بـ \code{sk} حقيقيّ؟ استخدم
  \code{StripeService}. وإلّا فالبديل \code{StubStripeService} لبيئة التطوير.
\item مِفتاح \code{Ai:Provider} يساوي \code{AzureOpenAi}؟ استخدم محوّل \en{Azure}.
  وإلّا فالمحوّل المحلّي \code{LocalAiService}.
\item مِفتاح \code{FileScanning:Enabled} مفعَّل؟ استخدم فحص \en{ClamAV} الحقيقيّ.
  وإلّا فمحوّل لا يفعل شيئًا.
\end{itemize}

\begin{why}
لماذا هذا أهمّ قرار في النظام؟ لسببين عمليّين ضخمين. \textbf{أولًا: القابلية
للاختبار.} نستطيع في الاختبارات استبدال \en{Stripe} و\en{Azure} و\en{ACS} ببدائل
حتميّة، فنختبر منطق العمل كلّه بلا إنترنت ولا تكلفة ولا بطء. \textbf{ثانيًا:
القابلية للتبديل.} نبدّل مزوّدًا كاملًا بتغيير سطر في الإعداد فقط — \code{Ai:Provider}،
\code{Storage:Provider}، \code{Email:Provider}، \code{FileScanning:Enabled}،
\code{Hangfire:Enabled} — بلا لمس كود المنطق. في التطوير كل شيء بدائل محلّية مجّانية،
وفي الإنتاج تُبدَّل بمزوّدات حقيقية عبر الإعداد. نفس الكود، بيئتان.
\end{why}

\begin{secnote}
هذا الفصل بين المنفذ والمحوّل هو أيضًا حصنٌ أمنيّ. المفاتيح السرّية والاتّصالات
الخارجية كلّها محبوسة في طبقة \en{Infrastructure} وحدها. طبقة العمل — قلب النظام —
لا تلمس مِفتاحًا سرّيًّا ولا تعرف عنوان خادم. لو أراد مُدقّق أمنيّ معرفة «أين تُستخدَم
المفاتيح؟»، فالإجابة في مكان واحد محصور، لا منثورة في كل مِلفّ.
\end{secnote}

\begin{qa}
«لماذا \en{Clean Architecture} بدل معماريّة طبقيّة بسيطة؟» — لأنّها تفرض قلب
الاعتماد \en{(Dependency Inversion)}: الدومين لا يعرف شيئًا عن \en{EF Core} أو
\en{Azure}، فنستطيع تبديل قاعدة البيانات أو أي خدمة دون لمس منطق العمل. أثبتنا هذا
عمليًّا بحقن بدائل \en{(stubs)} لكل الخدمات الخارجية في الاختبارات، وبتبديل مزوّد
الذكاء الاصطناعي عبر مِفتاح \code{Ai:Provider} وحده. في المعماريّة الطبقيّة
التقليدية تعتمد الطبقة العليا على الدنيا مباشرةً، فيصبح المنطق مرهونًا بالتفاصيل
التقنية — وهذا بالضبط ما نتجنّبه.
\end{qa}

\begin{summary}[الخلاصة في سطور]
\begin{itemize}
\item \en{ScholarPath} خلفيّتها مبنيّة على \en{Clean Architecture}: أربع طبقات
  متداخلة، والاعتماد يتّجه \textbf{للداخل فقط}.
\item الطبقات هي: \en{Domain} (القلب النقيّ) ثم \en{Application} (حالات الاستخدام)
  ثم \en{Infrastructure} (التقنيات) ثم \en{API} (الواجهة الرفيعة). أربعة مشاريع
  \en{.csproj} منفصلة، والترتيب مضمون بمراجع المشاريع.
\item \en{Domain} يحوي $\approx$55 كيانًا والـ \en{enums} وأحداث الدومين ومنفذين
  فقط \en{(}\code{ICurrentUserService}\en{, }\code{IDateTimeService}\en{)}، بلا أي
  اعتماد على \en{framework}.
\item \en{Application} يكتب حالات الاستخدام كـ \en{CQRS}، ومعها الـ \en{validators}
  و\en{AutoMapper}، ويُعلن $\approx$33 منفذًا يصف \emph{ماذا} يحتاج النظام.
\item \en{Infrastructure} يقدّم $\approx$46 محوّلًا تُنفّذ المنافذ بتقنيات حقيقية
  \en{(EF Core, Stripe, Azure OpenAI, ACS, Blob, Key Vault, SignalR, Hangfire)}.
\item القرار الأهمّ هو \en{Ports \& Adapters}: يجعل النظام \textbf{قابلًا للاختبار}
  و\textbf{قابلًا لتبديل المزوّد بالإعداد وحده} — \code{Ai:Provider}،
  \code{Storage:Provider}، \code{Email:Provider}، \code{FileScanning:Enabled}،
  \code{Hangfire:Enabled} — دون تغيير أي كود.
\end{itemize}
\end{summary}
